% ===================================================================
%  圖表第 2 號 — 人體。--- 小息。啲遊戲。
%
%  Traduction, faite en 2026, en cantonais ecrit d'aujourd'hui, en
%  caracteres traditionnels, sur l'ido de texto/io. Regles de la
%  colonne : voir 10-toubiu-01.tex.
%
%  LE CORPS HUMAIN EST L'ENDROIT OU LE CANTONAIS SE SEPARE LE PLUS DU
%  MANDARIN, ET CE N'EST PAS UN HASARD : ce sont les mots qu'on
%  apprend avant l'ecole, et l'ecole ne les a jamais remplaces.
%
%    脷    la langue      (Pekin : 舌)
%    膊頭  l'epaule       (Pekin : 肩膀)
%    手踭  le coude       (Pekin : 胳膊肘)
%    大髀  la cuisse      (Pekin : 大腿)
%    膝頭哥 le genou      (Pekin : 膝盖)
%    腳眼  la cheville    (Pekin : 脚踝)
%    腳踭  le talon       (Pekin : 脚后跟)
%    面珠  la joue        (Pekin : 脸颊)
%
%  Huit mots sur les quatre-vingts de l'anatomie, mais ce sont ceux
%  que le corps porte : la ou le tableau 1 se separait du mandarin par
%  la grammaire, celui-ci s'en separe par la chair.
%
%  ET LES CINQ SENS, EUX, SONT LES MEMES QU'A PEKIN — 視覺, 聽覺,
%  嗅覺, 味覺, 觸覺 — parce que ce sont des notions de manuel, non des
%  choses qu'on montre du doigt. C'est exactement le partage qu'on a
%  trouve en vietnamien.
%
%  LES JEUX SE PARTAGENT DE MEME, ET DEUX MOTS SUFFISENT A LE MONTRER :
%  « 波 » pour la balle, qui est l'anglais \textit{ball} passe par le
%  port de Hong Kong et que Pekin ne dit pas ; et « 單車 » pour la
%  bicyclette, la voiture a une place, la ou Pekin dit 自行车. On dit
%  aussi « 游水 », nager dans l'eau, et non 游泳.
% ===================================================================

%%K t02-apar-1 apar
\VUsaut{4.0mm}
\VUtitre{15.0pt}{60}{圖表第 2 號}
\VUsaut{2.0mm}
\VUtitre{11.4pt}{0}{人體。--- 小息。}
\VUtitre{11.4pt}{0}{啲遊戲。}

%%K t02-tit-0 sub
\VUtitre{13.2pt}{0}{人體。}
\VUsaut{2.4mm}
\VUcentre{10.2pt}{0}{\textit{（自然科學嘅簡單一課。）}}
\VUsaut{3.0mm}

%%K t02-tit-1 sub pos=t02-01-1
\VUpk{10.2pt}{40}{個頭。}
\VUsaut{3.0mm}

%%K t02-01-1 p
1. --- 人體嘅主要部分係：\VUgras{頭}、\VUgras{身}同\VUgras{四肢}。

%%K t02-02-1 p
2. --- 頭分兩部分：\VUgras{頭骨}，入面裝住\VUgras{腦}，外面畀\VUgras{頭髮}蓋住；
仲有\VUgras{塊面}。

%%K t02-03-1 p
3. --- 喺我哋幅圖上面，頭骨同腦都睇唔到；不過塊面重要嘅部分就睇得好清楚。

%%K t02-04-1 p
4. --- 首先係\VUgras{額頭}，佢顯出強大嘅才智同不斷運轉嘅思想。\VUgras{太陽穴}好
開闊。\VUgras{隻眼}靈活，張得好大。濃密嘅\VUgras{眉毛}同長長嘅\VUgras{眼睫毛}
護住佢。

%%K t02-05-1 p
5. --- 跟住係\VUgras{個鼻}同\VUgras{鼻孔}。個鼻幫我哋聞嘢同呼吸。鼻嘅兩邊係
\VUgras{面珠}，再遠啲就係\VUgras{對耳}。塊面下半部係\VUgras{把口}同\VUgras{下巴}。

%%K t02-06-1 p
6. --- 我哋睇得唔算好清楚，因為\VUgras{上唇鬚}同\VUgras{腮鬚}遮住咗佢哋。不過我哋
知道把口有兩片\VUgras{嘴唇}、\VUgras{上顎}同\VUgras{下顎}連埋佢哋嘅\VUgras{牙}，
仲有條\VUgras{脷}。我哋用把口講嘢同食嘢。我哋用條脷嚟試味。

%%K t02-07-1 p
7. --- 眼、耳、鼻同口，同啲手指一樣，都係五官嘅器官：\VUgras{視覺}、
\VUgras{聽覺}、\VUgras{嗅覺}、\VUgras{味覺}、\VUgras{觸覺}。

%%K t02-08-1 p
8. --- 所以個頭係人體最有趣嘅部分之一。

%%K t02-tit-2 sub
\VUsaut{4.4mm}
\VUpk{10.2pt}{40}{個身。}
\VUsaut{3.0mm}

%%K t02-09-1 p
9. --- \VUgras{條頸}接住個頭同個身。佢包括\VUgras{喉嚨}同\VUgras{頸背}。

%%K t02-10-1 p
10. --- 個身入面仲有好多重要嘅器官。\VUgras{胸腔}入面有\VUgras{肺}，我哋靠佢
呼吸；仲有\VUgras{個心}，佢係生命嘅中心。個心一停，生物就死。

%%K t02-11-1 p
11. --- \VUgras{肋骨}托住個胸腔。

%%K t02-12-1 p
12. --- 個身其餘嘅部分係\VUgras{腰}、\VUgras{個背}、\VUgras{膊頭}同\VUgras{個肚}。
我哋瞓覺嗰陣側住或者仰住瞓；我哋用膊頭孭嘢。

%%K t02-13-1 p
13. --- 個肚入面有\VUgras{個胃}同\VUgras{啲腸}。佢哋幫我哋消化食物。

%%K t02-tit-3 sub
\VUsaut{4.4mm}
\VUpk{10.2pt}{40}{四肢。}
\VUsaut{3.0mm}

%%K t02-14-1 p
14. --- 我哋有四條\VUgras{肢}：兩隻\VUgras{手臂}同兩隻\VUgras{腳}。我哋用手臂攞
嘢、揸嘢、抬嘢、執嘢；用腳企、行同跑。

%%K t02-15-1 p
15. --- \VUgras{膊頭}接住手臂同個身。一嚿好強壯嘅\VUgras{肌肉}，二頭肌，令手臂
郁到；\VUgras{手踭}就接住手臂同\VUgras{前臂}。\VUgras{手腕}組成\VUgras{隻手}嘅
關節，隻手末端係\VUgras{手指}同\VUgras{指甲}。喺隻手上面，我哋分得出\VUgras{靜脈}
（同動脈），\VUgras{血}就喺入面流。

%%K t02-16-1 p
16. --- 隻腳嘅部分係：\VUgras{大髀}、\VUgras{膝頭哥}同\VUgras{膝彎}、
\VUgras{小腿肚}——佢嘅\VUgras{肉}畀\VUgras{皮}包住——\VUgras{腳眼}同\VUgras{腳板}。
腳嘅\VUgras{骨}好粗好硬。腳前面係\VUgras{腳趾}。其餘嘅部分係\VUgras{腳底}同
\VUgras{腳踭}。

%%K t02-17-1 p
17. --- 呢個就係人體差唔多完整嘅描述。

%%K t02-17-2 p
\textit{註。} --- \textit{用「我……（個頭等等）痛」呢句講法，老師可以喺}
\VUgras{身體狀況}\textit{好定唔好呢個角度，將呢一整套詞語再用一次。}

%%K t02-tit-4 sub
\VUsaut{5.0mm}
\VUpk{10.2pt}{40}{體操。}
\VUsaut{2.4mm}
\VUcentre{10.2pt}{0}{\textit{（一位學生講嘅故事。）}}
\VUsaut{3.0mm}

%%K t02-18-1 p
18. --- 為咗身體柔軟、靈活同健康，我哋要練我哋嘅腳同手臂。所以我哋玩嘢，又做
\VUgras{體操}。

%%K t02-19-1 p
19. --- 平日呢兩樣都喺中學嘅操場度做（睇圖表第 1 號）。不過星期四同星期日，我哋
會去一塊大草地，喺嗰度有位跑同玩。

%%K t02-20-1 p
20. --- 我哋啲老師同父母有時同我哋一齊去；不過帶操嘅係\VUgras{體育老師}\textsuperscript{(12)}。

%%K t02-21-1 p
21. --- 草地上面，入口左邊有\VUgras{體操器械}\textsuperscript{(1)}；我哋爬
\VUgras{梯}\textsuperscript{(2)}，喺\VUgras{吊環}\textsuperscript{(3)} 同
\VUgras{吊槓}\textsuperscript{(4)} 上面翻筋斗，又用手攀上\VUgras{繩梯}\textsuperscript{(5)}
或者\VUgras{打結繩}\textsuperscript{(6)} 嘅頂。

%%K t02-22-1 p
22. --- 同時，我哋啲同學跳過\VUgras{木馬}\textsuperscript{(7)} 或者
\VUgras{雙槓}\textsuperscript{(11)}。第啲人就\VUgras{打拳}\textsuperscript{(8)} 或者
\VUgras{劍擊}\textsuperscript{(9)}，戴住面罩，攞住\VUgras{花劍}。最後仲有人
\VUgras{舉啞鈴}\textsuperscript{(10)}。

%%K t02-tit-5 sub
\VUsaut{4.6mm}
\VUpk{10.2pt}{40}{啲遊戲。}
\VUsaut{3.0mm}

%%K t02-23-1 p
23. --- 喺做體操嗰班人周圍，男仔女仔玩緊唔同嘅\VUgras{遊戲}。啲女仔玩
\VUgras{鐵環}\textsuperscript{(14)}；佢哋\VUgras{跳繩}\textsuperscript{(13)}，或者
叫自己嘅細佬同表兄弟一齊玩一局\VUgras{槌球}\textsuperscript{(15)}。佢哋最開心就係
喺對手以為好易穿過拱門嗰陣，用自己嘅\VUgras{木槌}撞開佢個\VUgras{波}！

%%K t02-24-1 p
24. --- 啲男仔中意用力嘅遊戲。佢哋鍾意打\VUgras{木瓶}\textsuperscript{(16)}，用一個
\VUgras{波}\textsuperscript{(17)} 好準咁撞冧佢哋。佢哋亦都唔嫌棄
\VUgras{陀螺}\textsuperscript{(18)} 同\VUgras{杯球}\textsuperscript{(19)}，甚至
\VUgras{青蛙戲}\textsuperscript{(20)}。不過佢哋最中意嘅仲係
\VUgras{波子}\textsuperscript{(21)} 同\VUgras{撞牆波}\textsuperscript{(22)}，因為
\VUgras{溫室}\textsuperscript{(26)} 嗰幅牆好啱畀輕波彈返轉頭。

%%K t02-25-1 p
25. --- 細個嘅約翰踩住\VUgras{高蹺}\textsuperscript{(24)}，好靈活，平衡好得好好。
喺佢附近，幾個同學喺溫室入面同\VUgras{籬笆}後面玩\VUgras{捉迷藏}\textsuperscript{(25)}。
第啲人就中意\VUgras{打手掌}\textsuperscript{(28)}，或者玩
\VUgras{羽毛球}\textsuperscript{(29)}，個球喺一支\VUgras{球拍}同另一支之間飛來飛去。

%%K t02-noto-1 noto
\VUnotes{143.2mm}{%
(*) 細路哥嘅遊戲成日有純粹本國嘅名。我哋盡量用伊多語為佢哋改名，或者靠佢哋本身
嘅動作，或者揀某國嘅名——只要嗰個名唔算太離譜。例如：\VUgras{木桶戲}（德文同法文）
就係\VUgras{青蛙戲} \textsuperscript{(20)}（英文），玩嘅人要將自己啲圓碟仔擲入嗰隻
張大口企喺穿窿箱（木桶）上面嘅鐵青蛙個口度。
\VUgras{跳膊頭}\textsuperscript{(30)} 同\VUgras{跳山羊}唔同嘅地方只有一樣：要跳過個
頭，玩嘅人要撐住同伴嘅膊頭，而喺「\VUgras{跳山羊}」入面就只係撐住佢個背。--- 又或者
有幾個人彎低身，第啲人撐住佢哋膊頭跳過佢哋個背；前者要頂得住撞力唔可以塌低，後者
要跳得過而唔失平衡（睇返呢張圖表）。%
}

%%K t02-26-1 p
26. --- 草地中間有兩棵樹。喺其中一棵樹腳下，七八個男仔組織咗一局
\VUgras{跳膊頭}\textsuperscript{(30)} (*)。呢個遊戲唔係個個國家都識；不過個個都知
\VUgras{跳山羊}係乜，雖然今次啲細路唔係玩緊呢個。

%%K t02-tit-6 sub
\VUsaut{5.0mm}
\VUpk{10.2pt}{40}{露天嘅遊戲。}
\VUsaut{3.4mm}

%%K t02-27-1 p
27. \VUgras{蒙眼捉人}\textsuperscript{(31)} 都好好玩；女仔同男仔輪流將
\VUgras{眼罩}戴上眼，去捉嗰啲唔夠醒目、畀人捉到嘅人。

%%K t02-28-1 p
28. --- 草地中間，呢個係好法國、不過有啲粗野嘅遊戲：叫做\VUgras{熊仔戲}\textsuperscript{(32)}，
佢比起靜靜哋嘅\VUgras{放紙鳶}\textsuperscript{(33)} 嘈好多，甚至比英國人嗰個
\VUgras{網球}\textsuperscript{(34)} 都嘈。

%%K t02-29-1 p
29. --- 最後嗰項運動係啲大學生嘅樂趣。連我哋啲老師都好樂意去打。佢要好多技巧同
敏捷，我好中意。至於\VUgras{鞦韆}\textsuperscript{(35)} 嘅樂趣，我就留返畀啲女仔。

%%K t02-30-1 p
30. --- 我唔中意英國人另一個遊戲，佢可能會變得危險。

%%K t02-30-1b p
我想講嘅係\VUgras{足球}\textsuperscript{(36)}，我大佬就好鍾意佢。我情願玩我哋舊時嗰個
\VUgras{奪旗戲}\textsuperscript{(38)}，一樣咁有益身體，而且真係冇咁粗野。

%%K t02-tit-7 sub
\VUsaut{4.6mm}
\VUpk{10.2pt}{40}{單車同球類。}
\VUsaut{3.0mm}

%%K t02-31-1 p
31. --- 我亦都好中意踩\VUgras{單車}\textsuperscript{(37)} 圍住草地兜一個圈。

%%K t02-32-1 p
32. --- 出年我會學\VUgras{騎馬}\textsuperscript{(39)}。一隻馬優雅咁行同小跑，
又飛奔跳過障礙，幾咁靚呀！

%%K t02-33-1 p
33. --- 夏天，我哋學\VUgras{游水}\textsuperscript{(40)}。我哋好開心咁潛落河、
喺清涼嘅河水入面浸。有時最後我哋會放一個紙\VUgras{氣球} \textsuperscript{(41)}，
佢一裝滿熱氣就威威咁升上天。

%%K t02-34-1 p
34. --- 仲有好多第啲遊戲，不過我數極都數唔完；睇完呢個簡略嘅講法，人人都明白，
逢星期四喺我哋塊靚草地度，係唔會悶嘅。

%%K t02-34-2 p
附註。--- \textit{老師可以好容易咁將每一個最重要嘅遊戲講得再詳細啲，去補足呢一章。}
\VUsaut{6.0mm}
\VUfilet{22mm}
